\chapter{2014年辽宁卷（文）}

\section{单选题}

\begin{problem}
已知全集 \(U = \R\)，\(A = \{x \mid x \leqslant 0\}\)，\(B = \{x \mid x \geqslant 1\}\)，则集合 \(\complement_{U}(A \cup B)=\)（\quad）

\choices
    {\( \{x \mid x \geqslant 0\} \)}
    {\(\{x\mid x\leqslant 1\}\)}
    {\(\{x\mid 0\leqslant x\leqslant 1\}\)}
    {\(\{x\mid 0 < x < 1\}\)}
\end{problem}

\begin{problem}
设复数 z 满足 \( (z - 2\i)(2 - 1) = 5 \) ，就是 z =（\quad）

\choices
    {\(2 + 3\i\)}
    {\(2 - 3\i\)}
    {\(3 + 2\i\)}
    {\(3 - 2\i\)}
\end{problem}

\begin{problem}
已知 \(a = 2^{-\frac{1}{2}}\)，\(b = \log_{\frac{1}{2}} \frac{1}{3}\)，\(c = \log_{\frac{1}{2}} \frac{1}{3}\)，则（\quad）

\choices
    {\(a > b > c\)}
    {\(a > c > b\)}
    {\(c > a > b\)}
    {\(c > b > a\)}
\end{problem}

\begin{problem}
已知 \(m, n\) 表示两条不同直线，\(\alpha\) 表示平面，下列说法正确的是（\quad）

\choices
    {若 \(m // \alpha, n // \alpha\)，则 \(m // n\)}
    {若 \(m \perp \alpha, n \subset \alpha\)，则 \(m \perp n\)}
    {若 \(m \perp \alpha, m \perp n\)，则 \(n // \alpha\)}
    {若 \(m // \alpha, m \perp n\)，则 \(n \perp \alpha\)}
\end{problem}

\begin{problem}
设 \(a, b, c\) 是非零向量，已知命题 \(p\) 若 \(a \cdot b = 0, b \cdot c = 0\)，则 \(a \cdot c = 0\)；命题 \(q\)；若 \(a // b, b // c\)，则 \(a // c\)，则下列命题中的真命题是（\quad）

\choices
    {\(p \vee q\)}
    {\(p \wedge q\)}
    {\((\neg p) \wedge \neg (q)\)}
    {\(p \vee (\neg q)\)}
\end{problem}

\begin{problem}
若将一个质点随机投入如图所示的长方形 \(ABCD\) 中，其中 \(AB = 2\)，\(BC = 1\)，则质点落在以 \(AB\) 为直径的半圆内的概率是（\quad）

\choices
    {\(\frac{\pi}{2}\)}
    {\(\frac{\pi}{4}\)}
    {\(\frac{\pi}{6}\)}
    {\(\frac{\pi}{8}\)}
\end{problem}

\begin{problem}
某几何体三视图如图所示，则该几何体的体积为（\quad）

\bitmapfigure[max width=0.56\linewidth,max totalheight=5.4cm]{img/2014/liaoning_liberal/q7_fig1.png}

\choices
    {\(8 - \frac{\pi}{4}\)}
    {\(8 - \frac{\pi}{2}\)}
    {\(8 - \pi\)}
    {\(8 - 2\pi\)}
\end{problem}

\begin{problem}
已知点 \(A(-2,3)\) 在抛物线 \(C: y^{2} = 2px\) 的准线上，记 \(C\) 的焦点为 \(F\)，则直线 \(AF\) 的斜率为（\quad）

\choices
    {\(-\frac{4}{3}\)}
    {\(-1\)}
    {\(-\frac{3}{4}\)}
    {\(-\frac{1}{2}\)}
\end{problem}

\begin{problem}
设等差数列 \(\{a_{n}\}\) 的公差为 \(d\)，若数列 \(\{2^{a_{1n - 1}}\}\) 为递减数列，则（\quad）

\choices
    {\(d < 0\)}
    {\(d > 0\)}
    {\(a_{1}d > 0\)}
    {\(a_{1}d < 0\)}
\end{problem}

\begin{problem}
已知 \( f(x) \) 为偶函数，当 \( x \geqslant 0 \) 时， \( f(x) = \begin{cases} \cos \pi x, & x \in \left[0, \frac{1}{2}\right], \\ 2x - 1, & x \in \left(\frac{1}{2}, +\infty\right). \end{cases} \) 则不等式 \(f(x - 1)\leqslant \frac{1}{2}\) 的解集为（\quad）

\choices
    {\(\left[\frac{1}{4}, \frac{2}{3}\right] \cup \left[\frac{4}{3}, \frac{7}{4}\right]\)}
    {\(\left[-\frac{3}{4}, -\frac{1}{3}\right] \cup \left[\frac{1}{4}, \frac{2}{3}\right]\)}
    {\(\left[\frac{1}{3}, \frac{3}{4}\right] \cup \left[\frac{4}{3}, \frac{7}{4}\right]\)}
    {\(\left[-\frac{3}{4}, -\frac{1}{3}\right] \cup \left[\frac{1}{3}, \frac{3}{4}\right]\)}
\end{problem}

\begin{problem}
将函数 \( y = 3\sin \left(2x + \frac{\pi}{3}\right) \) 的图象向右平移 \( \frac{\pi}{2} \) 个单位长度，所得图象对应的函数（\quad）

\choices
    {在区间 \(\left[\frac{\pi}{12}, \frac{7\pi}{12}\right]\) 上单调递减}
    {在区间 \(\left[\frac{\pi}{12}, \frac{7\pi}{12}\right]\) 上单调递增}
    {在区间 \(\left[-\frac{\pi}{6}, \frac{\pi}{3}\right]\) 上单调递减}
    {在区间 \(\left[-\frac{\pi}{6}, \frac{\pi}{3}\right]\) 上单调递增}
\end{problem}

\begin{problem}
当 \(x \in [-2, 1]\) 时，不等式 \(ax^3 - x^2 + 4x + 3 \geqslant 0\) 恒成立，则实数 \(a\) 的取值范围是（\quad）

\choices
    {\([-5, - 3]\)}
    {\(\left[-6, -\frac{9}{8}\right]\)}
    {\([-6, - 2]\)}
    {\([-4, - 3]\)}
\end{problem}

\section{填空题}

\begin{problem}
执行如图的程序框图，若输入 n = 3，则输出 T =\fillinblank{} \fillinblank{}.
\end{problem}

\begin{problem}
已知 \(x, y\) 满足条件 \(\left\{ \begin{array}{l} 2x + y - 2 \geqslant 0, \\ x - 2y + 4 \geqslant 0, \\ 3x - y - 3 \leqslant 0. \end{array} \right.\) 则目标函数 \(z = 3x + 4y\) 的最大值为 \fillinblank{} \fillinblank{}.
\end{problem}

\begin{problem}
已知椭圆 \(C: \frac{x^2}{9} + \frac{y^2}{4} = 1\)，点 \(M\) 与 \(C\) 的焦点不重合，若 \(M\) 关于 \(C\) 的焦点的对称点分别为 \(A, B\)，线段 \(MN\) 的中点在 \(C\) 上，则 \(|AN| + |BN| = \fillinblank{}\). 16. 对于 \(c > 0\)，当非零实数 \(a, b\) 满足 \(4a^2 - 2ab + b^2 - c = 0\) 且使 \(|2a + b|\) 最大时，\(\frac{a}{a + b} + \frac{b}{b + a} + \frac{1}{c}\) 的最小值为 \(\fillinblank{}\). 
\end{problem}

\section{解答题}

\begin{problem}
在 \(\triangle ABC\) 中，内角 \(A, B, C\) 的对边分别为 \(a, b, c\)，且 \(a > c\)，已知 \(\overrightarrow{BA} \cdot \overrightarrow{BC} = 2\)，\(c = 08\)，\(b = \frac{1}{3}\)，\(b = 3\). 求: (1) \(a\) 和 \(c\) 的值; (2) \( \cos(B-C) \) 的值.
\end{problem}

\begin{problem}
某大学餐饮中心为了了解新生的饮食习惯，在全校一年级学生中进行了抽样调查，调查结果如下表所示： 喜欢甜品 不喜欢甜品 合计 南方学生 60 20 80 北方学生 10 10 20 合计 70 30 100 (1) 根据表中数据，问是否有 95\% 的把握认为“南方学生和北方学生在选用甜品的饮食习惯方面有差异”; (2) 已知在被调查的北方学生中有 5 名数学系的学生，其中 2 名喜欢甜品，现在从这 5 名学生中随机抽取 3 人，要求至少有 1 人喜欢甜品的概率. 附： \(\chi^2 = \frac{n(n_1n_2 - n_1n_2n_2)^2}{n_1 + n_2 + n_1 + n_2x}\left| \begin{array}{llll} P(x^2\geqslant k) & 0.100 & 0.050 & 0.010\\ k & 2.706 & 3.841 & 6.635 \end{array} \right|\)
\end{problem}

\begin{problem}
如图，\(\triangle ABC\) 和 \(\triangle BCD\) 所在平面互相垂直，且 \(AB = BC = BD = 2\)，\(\angle ABC = \angle DBC = 120^{\circ}\)，\(E, F, G\) 分别为 \(AC, DC, AD\) 的中点. (1) 求证: \( EF \perp \) 平面 BCG; (2) 求三棱锥 D - BCC 的体积. 附: 锥体的体积公式 \( V = \frac{3}{2} Sh_{1} \) ，其中 S 为底面面积，h 为高.
\end{problem}

\begin{problem}
已知函数 \( f(x) = \pi (x - \cos x) - 2\sin x - 2, g(x) = (x - \pi)\sqrt{\frac{1 - \sin x}{1 + \sin x}} + \frac{2\pi}{\pi} - 1. \) (1) 证明: 存在唯一 \(x_0 \in \left(0, \frac{\pi}{2}\right)\)，使 \(f(x_0) = 0\); (2) 证明: 存在唯一 \(x_{1} \in \left(\frac{\pi}{2}, \pi\right)\)，使 \(g(x_{1}) = 0\)，且对 (1) 中的 \(x_{0}\)，有 \(x_{0} + x_{1} > \pi\).
\end{problem}

\begin{problem}
将圆 \(x^{2} + y^{2} = 1\) 上每一点的横坐标保持不变，纵坐标变为原来的 2 倍，得曲线 \(C\). (1) 写出 C 的参数方程; (2) 设直线 \(l: 2x + y - 2 = 0\) 与 \(C\) 的交点为 \(P_{1}, P_{2}\)，以坐标原点为极点，\(x\) 轴正半轴为极轴建立极坐标系. 求过线段 \(P_{1}P_{2}\) 的中点且与 \(l\) 垂直的直线的极坐标方程.
\end{problem}

\begin{problem}
如图，\(EP\) 交圆于 \(E, C\) 两点，\(PD\) 切圆于 \(D, G\) 为 \(CE\) 上一点且 \(PG = PD\)，连接 \(DG\) 并延长交圆于点 \(A\)，作弦 \(AB\) 垂直 \(EP\)，垂足为 \(F\). (1) 求证: AB 为圆的直径; (2) 若 AC = BD，求证: AB = ED.
\end{problem}

\begin{problem}
设函数 \( f(x) = 2|x - 1| + x - 1 \)，\( g(x) = 16x^{2} - 8x + 1 \)，记 \( f(x) \leqslant 1 \) 的解集为 \( M \)，\( g(x) \leqslant 4 \) 的解集为 \( N \). (1) 求 M; (2) 当 \(x \in M \cap N\) 时，证明: \(x^2 f(x) + x[f(x)]^9 \leqslant \frac{1}{4}\).
\end{problem}

